%% kapa: exam=1 round=37 year=2026 period=1 form=A subject=real_estate subject_ko=부동산학원론 seq=2 number=82
%% source: Q-Net 감정평가사 기출문제 (cst00309) artlSeq=5258924 file=raw/1008/5258924/2026년 제37회 감정평가사 1차 1교시.pdf
%% conversion: kordoc 4.24.51 -> tools/gen1.py (auto); review=proofread-sonnet; flags=none
\begin{question}{82}
부동산과 준부동산에 관한 설명으로 옳지 않은 것은? (다툼이 있으면 판례에 따름)
\begin{choices}
  \item 적재용량 25톤인 덤프트럭은 등록을 통해 소유권을 공시한다.
  \item 총톤수 25톤인 범선(帆船)은 등기를 통해 소유권을 공시한다.
  \item 토지에서 경작되는 농작물은 명인방법을 갖추더라도 토지의 일부로 취급된다.
  \item 토지에 부착된 한 그루의 수목은 「입목에 관한 법률」에 따른 입목등기를 할 수 없다.
  \item 토지에 부착된 한 그루의 수목은 명인방법을 갖추면 토지와는 독립한 물건으로 취급된다.
\end{choices}
\end{question}
\answer{3}
